% Generated draft: requires mathematical and source review.
\documentclass[11pt]{article}
\usepackage[margin=25mm]{geometry}
\usepackage{fontspec}
\setmainfont{Latin Modern Roman}
\usepackage{amsmath,amssymb,mathtools,unicode-math}
\setmathfont{Latin Modern Math}
\usepackage{enumitem,needspace,xurl,hyperref}
\setlength{\parindent}{0pt}
\setlength{\parskip}{5pt}
\setlength{\emergencystretch}{3em}
\newcommand{\N}{\mathbb{N}}
\newcommand{\Z}{\mathbb{Z}}
\newcommand{\Q}{\mathbb{Q}}
\newcommand{\R}{\mathbb{R}}
\newcommand{\C}{\mathbb{C}}
\newcommand{\F}{\mathbb{F}}
\newcommand{\A}{\mathbb{A}}
\newcommand{\PP}{\mathbb{P}}
\newcommand{\E}{\mathbb{E}}
\newcommand{\eps}{\varepsilon}
\newcommand{\dd}{\,\mathrm{d}}
\begin{document}
\section*{UnsolvedMath ID: 30001799}
\textbf{Source problem code:} catalogue entry
\subsection{Definitions and mathematical statement}
Write \(K_m(R)\) for the algebraic \(K\)-group of a ring \(R\) in degree \(m\), and write \(F_n^\gamma K_m(R)\) for the subgroup denoted by the \(n\)-th step of the gamma filtration in the supplied record. The record does not specify the ring conventions, the coefficient convention, the construction of this filtration on higher \(K\)-groups, or the allowed filtration indices.

The supplied formulation of the Beilinson--Soul\'e vanishing conjecture asks whether, for every regular ring \(R\) and every admissible filtration index \(n\),
\[
K_m(R)=F_n^\gamma K_m(R)
\qquad\text{for every integer }m\geq\max\{2n-1,1\}.
\]
Here ``admissible filtration index'' retains an unresolved source qualification: the record does not give a domain for \(n\). Likewise, ``regular ring'' is retained without adding unstated commutativity, Noetherianity, or dimension hypotheses.

The record also describes this as equivalent to vanishing of Adams eigenspaces \(K_m(R)^{(n)}\) in an ``excluded range,'' except for \(K_0(R)^{(0)}\). Under the usual rational interpretation, when Adams operations \(\psi^k\) are defined, the notation would mean
\[
K_m(R)^{(n)}
=
\bigl\{x\in K_m(R)\otimes_{\Z}\Q:
\psi^k(x)=k^n x\text{ for every integer }k\geq1\bigr\}.
\]
This rational interpretation is explanatory and is not specified by the record. The record supplies neither explicit inequalities defining the excluded range nor the hypotheses and filtration conventions needed to establish the asserted equivalence. In particular, the boundary degrees of the eigenspace formulation remain unresolved. The exception \(K_0(R)^{(0)}\) is preserved as supplied; it does not assert that this group vanishes.

The accompanying literature assessment says that the entry combines Beilinson--Soul\'e and Parshin vanishing statements, but no separate Parshin statement or its hypotheses is supplied. Consequently, no additional Parshin assertion is formulated here.
\subsection{Short English statement}
The supplied question asks whether sufficiently high algebraic \(K\)-groups of every regular ring lie entirely in a specified gamma-filtration step. Its claimed equivalent Adams-eigenspace formulation needs clarification of coefficients, index ranges, boundary degrees, and ring hypotheses. The supplied status is ``partially solved''; this is a dated curation label, not an independently verified mathematical conclusion.
\subsection{Sources}
\begin{itemize}
\item Local UnsolvedMath registry: supplied complete record for permanent ID \(30001799\), source problem code \texttt{catalogue entry}, titled ``Beilinson--Soul\'e Vanishing Conjecture.'' The record belongs to the Oberwolfach Reports Open Problems collection and gives proposed year 2011; no individual proposer is identified.
\item Supplied report citation: \emph{New Directions in Algebraic \(K\)-Theory} (2011), Oberwolfach Reports, pp.~1469--1509, DOI \texttt{10.4171/owr/2011/27}. Supplied source URL and literature-source link labelled ``original survey'': \url{https://doi.org/10.4171/owr/2011/27}. The background describes Oberwolfach Reports as documenting mathematical workshops and their open-problem discussions.
\item Supplied literature-source link labelled ``rational Parshin theorem'': \url{https://arxiv.org/abs/2011.05544}. The record supplies no authors, title, theorem statement, or explanation connecting this reference to the displayed filtration assertion.
\item Literature assessment dated exactly \texttt{2026-08-21}: the supplied curation labels the entry ``partially solved'' and says that it combines Beilinson--Soul\'e and Parshin vanishing statements, that low dimensional and rational cases are known, and that the general assertions remain open. These are supplied dated curation claims, not independently verified here. The background explicitly requests independent verification of time-sensitive status and cited claims.
\item The registry labels the statement \texttt{corrected\_verified} and reports a manual semantic, typographic, and LaTeX review of a meaning-preserving cleanup. That supplied review description does not resolve the mathematical formulation issues recorded here or establish the conjecture.
\end{itemize}
\end{document}
